\chapter{Hypothesis testing}
\label{ch:hypothesis-testing}

\begin{quote}
\emph{``The purpose of a statistical test is not to prove something is true. It is to measure how surprised we should be if nothing is going on.''}
\end{quote}

% ============================================================
\section{The logic of hypothesis testing}

Every clinical question can be framed as a test between two competing statements:

\begin{itemize}
  \item \textbf{Null hypothesis ($H_0$):} There is \emph{no} effect, no difference, no association. The drug does not work. The two groups have the same blood pressure.
  \item \textbf{Alternative hypothesis ($H_1$):} There \emph{is} an effect, a difference, or an association. The drug lowers cholesterol. The treatment group has lower blood pressure than the control.
\end{itemize}

We collect data, compute a test statistic, and ask: \emph{if $H_0$ were true, how likely is it that we would see data this extreme?} That probability is the \textbf{p-value}.

\begin{clinicalnote}
Clinical framing matters. A statistician writes $H_0: \mu_1 = \mu_2$. A clinician writes ``there is no difference in systolic blood pressure between the statin group and the placebo group.'' Both say the same thing --- use whichever makes the research question clear.
\end{clinicalnote}

\subsection{The steps of a hypothesis test}

\begin{enumerate}
  \item \textbf{State the hypotheses} ($H_0$ and $H_1$).
  \item \textbf{Choose the significance level} $\alpha$ (usually 0.05).
  \item \textbf{Collect data} and compute the \textbf{test statistic}.
  \item \textbf{Compute the p-value} --- the probability of observing a result at least as extreme as yours, assuming $H_0$ is true.
  \item \textbf{Decide:} if $p < \alpha$, reject $H_0$. Otherwise, fail to reject $H_0$.
\end{enumerate}

\begin{minted}{python}
import numpy as np
import pandas as pd
from scipy import stats
import matplotlib.pyplot as plt
import seaborn as sns
\end{minted}

% ============================================================
\section{p-values: what they mean and what they don't}

\subsection{What a p-value IS}

A p-value is the probability of observing data as extreme as (or more extreme than) what you collected, \emph{assuming} $H_0$ is true. A $p = 0.03$ means: if the drug truly has no effect, there is a 3\% chance of seeing results this extreme by random chance alone.

\subsection{What a p-value is NOT}

\begin{itemize}
  \item It is \textbf{not} the probability that $H_0$ is true.
  \item It is \textbf{not} the probability that the result is due to chance.
  \item It does \textbf{not} measure the size of the effect.
  \item $p = 0.049$ is not meaningfully different from $p = 0.051$.
\end{itemize}

\begin{warningbox}
A small p-value does not mean a \emph{large} or \emph{clinically important} effect. A study with 100\,000 patients can find $p < 0.001$ for a blood pressure difference of 0.5~mmHg --- statistically significant but clinically meaningless. Always report \textbf{effect size} alongside p-values.
\end{warningbox}

\subsection{Type I and Type II errors}

\begin{itemize}
  \item \textbf{Type I error (false positive):} You reject $H_0$ when it is actually true. You conclude the drug works when it does not. Controlled by $\alpha$.
  \item \textbf{Type II error (false negative):} You fail to reject $H_0$ when $H_1$ is actually true. You miss a real effect. Related to \textbf{power} ($1 - \beta$).
\end{itemize}

\begin{minted}{python}
# Visualize the logic: sampling distribution under H0
np.random.seed(42)
null_distribution = np.random.normal(loc=0, scale=1, size=10000)

fig, ax = plt.subplots(figsize=(9, 4))
ax.hist(null_distribution, bins=60, density=True, alpha=0.7, color="steelblue")
ax.axvline(1.96, color="red", linestyle="--", label="Critical value (alpha=0.05)")
ax.axvline(-1.96, color="red", linestyle="--")
ax.fill_betweenx([0, 0.45], 1.96, 3.5, alpha=0.3, color="red", label="Rejection region")
ax.fill_betweenx([0, 0.45], -3.5, -1.96, alpha=0.3, color="red")
ax.set_xlabel("Test statistic (z)")
ax.set_ylabel("Density")
ax.set_title("Two-tailed test: rejection regions under H0")
ax.legend()
plt.tight_layout()
plt.show()
\end{minted}

% ============================================================
\section{One-sample t-test}

\textbf{Question:} Is the mean total cholesterol in our clinic patients different from the national average of 200~mg/dL?

\begin{minted}{python}
# Simulated clinic cholesterol data (mg/dL)
np.random.seed(42)
cholesterol = np.random.normal(loc=212, scale=35, size=80)

# One-sample t-test: is the mean different from 200?
t_stat, p_value = stats.ttest_1samp(cholesterol, popmean=200)
print(f"Sample mean: {cholesterol.mean():.1f} mg/dL")
print(f"t-statistic: {t_stat:.3f}")
print(f"p-value: {p_value:.4f}")

if p_value < 0.05:
    print("Reject H0: mean cholesterol differs from 200 mg/dL")
else:
    print("Fail to reject H0: no evidence mean differs from 200 mg/dL")
\end{minted}

\begin{clinicalnote}
The one-sample t-test compares a sample mean to a known reference value. In clinical practice, the reference might be a national average, a guideline threshold, or a pre-treatment baseline.
\end{clinicalnote}

% ============================================================
\section{Two-sample t-test}

\textbf{Question:} Is systolic blood pressure different between the treatment group and the control group?

\begin{minted}{python}
# Simulated clinical trial data
np.random.seed(42)
treatment = np.random.normal(loc=128, scale=15, size=60)
control = np.random.normal(loc=138, scale=16, size=55)

# Independent two-sample t-test
t_stat, p_value = stats.ttest_ind(treatment, control)
print(f"Treatment mean: {treatment.mean():.1f} mmHg")
print(f"Control mean:   {control.mean():.1f} mmHg")
print(f"Difference:     {control.mean() - treatment.mean():.1f} mmHg")
print(f"t-statistic:    {t_stat:.3f}")
print(f"p-value:        {p_value:.4f}")
\end{minted}

\subsection{Checking assumptions}

The t-test assumes approximate normality and equal variances. Check both:

\begin{minted}{python}
# Normality: Shapiro-Wilk test (H0: data is normally distributed)
_, p_treat = stats.shapiro(treatment)
_, p_ctrl = stats.shapiro(control)
print(f"Shapiro-Wilk p (treatment): {p_treat:.4f}")
print(f"Shapiro-Wilk p (control):   {p_ctrl:.4f}")

# Equal variances: Levene's test
_, p_levene = stats.levene(treatment, control)
print(f"Levene's test p: {p_levene:.4f}")

# If variances are unequal, use Welch's t-test:
t_stat, p_value = stats.ttest_ind(treatment, control, equal_var=False)
print(f"Welch's t-test p-value: {p_value:.4f}")
\end{minted}

\begin{pythontip}
Set \texttt{equal\_var=False} in \texttt{ttest\_ind()} to use Welch's t-test. It does not assume equal variances and is often the safer default.
\end{pythontip}

% ============================================================
\section{Paired t-test}

\textbf{Question:} Does a 12-week exercise intervention reduce systolic blood pressure?

Each patient is measured \emph{before} and \emph{after}. The paired t-test uses the within-patient differences.

\begin{minted}{python}
# Simulated before/after intervention data
np.random.seed(42)
n_patients = 45
bp_before = np.random.normal(loc=142, scale=12, size=n_patients)
# After intervention: mean reduction of ~8 mmHg with individual variation
bp_after = bp_before - np.random.normal(loc=8, scale=6, size=n_patients)

# Paired t-test
t_stat, p_value = stats.ttest_rel(bp_before, bp_after)
differences = bp_before - bp_after
print(f"Mean BP before:     {bp_before.mean():.1f} mmHg")
print(f"Mean BP after:      {bp_after.mean():.1f} mmHg")
print(f"Mean reduction:     {differences.mean():.1f} mmHg")
print(f"t-statistic:        {t_stat:.3f}")
print(f"p-value:            {p_value:.6f}")
\end{minted}

\begin{minted}{python}
# Visualize before/after
fig, axes = plt.subplots(1, 2, figsize=(12, 5))

# Paired lines
for i in range(n_patients):
    axes[0].plot([0, 1], [bp_before[i], bp_after[i]],
                 color="gray", alpha=0.4)
axes[0].plot([0, 1], [bp_before.mean(), bp_after.mean()],
             color="red", linewidth=3, marker="o", markersize=8)
axes[0].set_xticks([0, 1])
axes[0].set_xticklabels(["Before", "After"])
axes[0].set_ylabel("Systolic BP (mmHg)")
axes[0].set_title("Individual patient trajectories")

# Distribution of differences
axes[1].hist(differences, bins=15, edgecolor="black", alpha=0.7, color="steelblue")
axes[1].axvline(0, color="red", linestyle="--", label="No change")
axes[1].axvline(differences.mean(), color="green", linestyle="--",
                label=f"Mean = {differences.mean():.1f}")
axes[1].set_xlabel("BP reduction (mmHg)")
axes[1].set_ylabel("Frequency")
axes[1].set_title("Distribution of BP changes")
axes[1].legend()

plt.tight_layout()
plt.show()
\end{minted}

\begin{warningbox}
A common mistake is using an \emph{independent} two-sample t-test on before/after data. This ignores the pairing and loses statistical power. If the same patients are measured twice, always use a \textbf{paired} t-test.
\end{warningbox}

% ============================================================
\section{Chi-square test of independence}

\textbf{Question:} Is diabetes prevalence associated with obesity category?

The chi-square test works with \emph{categorical} variables. It compares observed counts to what we would expect if the variables were independent.

\begin{minted}{python}
# Simulated cross-tabulation data
np.random.seed(42)
n = 500
bmi_cat = np.random.choice(["Normal", "Overweight", "Obese"],
                           size=n, p=[0.3, 0.4, 0.3])
# Diabetes probability increases with obesity
diabetes_prob = {"Normal": 0.08, "Overweight": 0.18, "Obese": 0.35}
diabetes = [np.random.binomial(1, diabetes_prob[b]) for b in bmi_cat]

df_chi = pd.DataFrame({"bmi_category": bmi_cat, "diabetes": diabetes})
df_chi["diabetes_label"] = df_chi["diabetes"].map({0: "No", 1: "Yes"})

# Contingency table
contingency = pd.crosstab(df_chi["bmi_category"], df_chi["diabetes_label"])
print(contingency)
print()

# Chi-square test
chi2, p_value, dof, expected = stats.chi2_contingency(contingency)
print(f"Chi-square statistic: {chi2:.2f}")
print(f"Degrees of freedom:   {dof}")
print(f"p-value:              {p_value:.4f}")
print(f"\nExpected frequencies (if independent):")
print(pd.DataFrame(expected, index=contingency.index,
                   columns=contingency.columns).round(1))
\end{minted}

\begin{clinicalnote}
The chi-square test requires that expected cell counts are at least 5. If you have small cells (e.g., a rare disease), use Fisher's exact test instead: \texttt{stats.fisher\_exact(table)} for $2 \times 2$ tables.
\end{clinicalnote}

% ============================================================
\section{Mann-Whitney U test}

When your data is not normally distributed --- common with hospital length-of-stay, cost data, or small samples --- use the Mann-Whitney U test. It compares the \emph{ranks} of observations rather than the means.

\begin{minted}{python}
# Hospital length of stay: surgical vs non-surgical (right-skewed data)
np.random.seed(42)
los_surgical = np.random.exponential(scale=7, size=40)
los_medical = np.random.exponential(scale=4.5, size=45)

# Mann-Whitney U test
u_stat, p_value = stats.mannwhitneyu(los_surgical, los_medical,
                                      alternative="two-sided")
print(f"Median LOS (surgical): {np.median(los_surgical):.1f} days")
print(f"Median LOS (medical):  {np.median(los_medical):.1f} days")
print(f"U statistic:           {u_stat:.0f}")
print(f"p-value:               {p_value:.4f}")
\end{minted}

\begin{minted}{python}
# Visualize the skewed distributions
fig, ax = plt.subplots(figsize=(8, 4))
ax.hist(los_surgical, bins=15, alpha=0.6, label="Surgical", color="coral")
ax.hist(los_medical, bins=15, alpha=0.6, label="Medical", color="steelblue")
ax.set_xlabel("Length of stay (days)")
ax.set_ylabel("Frequency")
ax.set_title("Hospital length of stay by department")
ax.legend()
plt.tight_layout()
plt.show()
\end{minted}

\begin{pythontip}
Rule of thumb for choosing a test:
\begin{itemize}
  \item Normal data, two groups $\rightarrow$ t-test
  \item Non-normal data, two groups $\rightarrow$ Mann-Whitney U
  \item Categorical data $\rightarrow$ chi-square
  \item Paired measurements $\rightarrow$ paired t-test (or Wilcoxon signed-rank if non-normal)
\end{itemize}
\end{pythontip}

% ============================================================
\section{Multiple testing correction}

\subsection{The problem}

If you test 20 hypotheses at $\alpha = 0.05$, you expect $20 \times 0.05 = 1$ false positive \emph{even if nothing is going on}. In genomics, you might test 20\,000 genes simultaneously. Without correction, you would get 1\,000 spurious ``discoveries.''

\begin{minted}{python}
# Simulation: 1000 tests where H0 is true for all
np.random.seed(42)
n_tests = 1000
p_values = []
for _ in range(n_tests):
    # Two random samples from the SAME distribution (H0 is true)
    a = np.random.normal(0, 1, 30)
    b = np.random.normal(0, 1, 30)
    _, p = stats.ttest_ind(a, b)
    p_values.append(p)

p_values = np.array(p_values)
false_positives = (p_values < 0.05).sum()
print(f"Tests performed: {n_tests}")
print(f"False positives (p < 0.05): {false_positives}")
print(f"False positive rate: {false_positives/n_tests:.1%}")
\end{minted}

\subsection{Bonferroni correction}

The simplest fix: divide $\alpha$ by the number of tests.

\begin{minted}{python}
from statsmodels.stats.multitest import multipletests

# Bonferroni correction
rejected_bonf, pvals_corrected_bonf, _, _ = multipletests(
    p_values, alpha=0.05, method="bonferroni"
)
print(f"Bonferroni: {rejected_bonf.sum()} significant results")
\end{minted}

\subsection{Benjamini-Hochberg (FDR)}

Less conservative: controls the \emph{false discovery rate} (the proportion of rejected hypotheses that are false positives).

\begin{minted}{python}
# Benjamini-Hochberg FDR correction
rejected_fdr, pvals_corrected_fdr, _, _ = multipletests(
    p_values, alpha=0.05, method="fdr_bh"
)
print(f"FDR (BH): {rejected_fdr.sum()} significant results")
\end{minted}

\begin{clinicalnote}
In genomics and proteomics, FDR correction is standard. When you read a paper reporting ``542 differentially expressed genes at FDR $< 0.05$,'' the authors used Benjamini-Hochberg (or a similar method) to control for the thousands of simultaneous tests.
\end{clinicalnote}

\begin{warningbox}
Bonferroni is very conservative --- it minimizes false positives but can miss real effects (high Type~II error). FDR is a good middle ground for exploratory analyses. For confirmatory clinical trials, Bonferroni is often preferred because false positives are more costly.
\end{warningbox}

% ============================================================
\section{Effect size: clinical vs statistical significance}

\subsection{Cohen's d}

Cohen's $d$ measures the size of the difference in \emph{standard deviation units}:
$$d = \frac{\bar{x}_1 - \bar{x}_2}{s_{\text{pooled}}}$$

\begin{itemize}
  \item $|d| < 0.2$: negligible effect
  \item $0.2 \leq |d| < 0.5$: small effect
  \item $0.5 \leq |d| < 0.8$: medium effect
  \item $|d| \geq 0.8$: large effect
\end{itemize}

\begin{minted}{python}
def cohens_d(group1, group2):
    """Compute Cohen's d for two independent groups."""
    n1, n2 = len(group1), len(group2)
    var1, var2 = group1.var(ddof=1), group2.var(ddof=1)
    pooled_std = np.sqrt(((n1 - 1) * var1 + (n2 - 1) * var2) / (n1 + n2 - 2))
    return (group1.mean() - group2.mean()) / pooled_std

# Using our treatment vs control BP data
d = cohens_d(treatment, control)
print(f"Cohen's d: {d:.2f}")
print(f"Interpretation: {'large' if abs(d) >= 0.8 else 'medium' if abs(d) >= 0.5 else 'small' if abs(d) >= 0.2 else 'negligible'} effect")
\end{minted}

\subsection{Why effect size matters}

\begin{minted}{python}
# Large sample + tiny effect = significant p-value
np.random.seed(42)
huge_group1 = np.random.normal(120.0, 15, size=50000)
huge_group2 = np.random.normal(120.5, 15, size=50000)  # only 0.5 mmHg difference

t, p = stats.ttest_ind(huge_group1, huge_group2)
d = cohens_d(huge_group1, huge_group2)
print(f"Mean difference: {huge_group2.mean() - huge_group1.mean():.2f} mmHg")
print(f"p-value: {p:.6f}")
print(f"Cohen's d: {d:.3f}")
print("Statistically significant? Yes. Clinically meaningful? No.")
\end{minted}

\begin{clinicalnote}
In a clinical trial report, the FDA and EMA expect both p-values \emph{and} confidence intervals for the treatment effect. A drug that lowers HbA1c by 0.01\% with $p < 0.001$ will not get approved. The effect must be \textbf{clinically meaningful}.
\end{clinicalnote}

% ============================================================
\section{Practical: hypothesis testing with Framingham-style data}

\begin{minted}{python}
# Load a Framingham-style cardiovascular dataset
url = ("https://raw.githubusercontent.com/dsrscientist/"
       "dataset-for-machine-learning/master/framingham.csv")
fram = pd.read_csv(url)
print(f"Shape: {fram.shape}")
print(f"Columns: {fram.columns.tolist()}")
fram.head()
\end{minted}

\begin{minted}{python}
# Quick exploration
fram.describe()
\end{minted}

\begin{minted}{python}
# Drop rows with missing values for simplicity
fram_clean = fram.dropna()
print(f"Clean dataset: {fram_clean.shape[0]} patients")
\end{minted}

\begin{exercise}
Using the Framingham dataset:
\begin{enumerate}
  \item \textbf{One-sample t-test.} Is the mean total cholesterol in this cohort different from the US national average of 200~mg/dL?
  \item \textbf{Two-sample t-test.} Compare systolic blood pressure between patients who developed CHD (\texttt{TenYearCHD = 1}) and those who did not. Report the p-value, mean difference, and Cohen's $d$.
  \item \textbf{Chi-square test.} Create a $2 \times 2$ contingency table of \texttt{currentSmoker} vs \texttt{TenYearCHD}. Is smoking associated with 10-year CHD risk?
  \item \textbf{Mann-Whitney U.} Compare \texttt{cigsPerDay} between CHD and non-CHD groups. Why is Mann-Whitney more appropriate than a t-test here?
  \item \textbf{Multiple testing.} Run t-tests comparing CHD vs non-CHD for all continuous variables (age, totChol, sysBP, diaBP, BMI, heartRate, glucose). Apply Bonferroni and FDR correction. Which variables remain significant after each correction?
  \item \textbf{Clinical interpretation.} For the variable with the largest effect size (Cohen's $d$), discuss whether the difference is clinically meaningful.
\end{enumerate}
\end{exercise}

% ============================================================
\section{Chapter summary}

\begin{itemize}
  \item A hypothesis test measures how surprising the data would be if $H_0$ were true.
  \item The p-value is \emph{not} the probability that $H_0$ is true.
  \item One-sample t-test: compare a sample mean to a reference value.
  \item Two-sample t-test: compare means between two independent groups.
  \item Paired t-test: compare means from the same subjects measured twice.
  \item Chi-square test: test association between two categorical variables.
  \item Mann-Whitney U: non-parametric alternative when data is not normal.
  \item Multiple testing correction (Bonferroni, FDR) prevents false discoveries.
  \item Effect size (Cohen's $d$) distinguishes statistical from clinical significance.
  \item Always report both p-values and effect sizes in clinical research.
\end{itemize}
